\chapter{Optimisations and Extensions}\label{ch:opt}

The first evaluation of the Phase~I scheme showed where its cost lies. Every BSM carried a full
ticket and a signature of about 330 bytes, sent ten times a second by every vehicle. At 150 vehicles
the safety traffic alone occupied about a quarter of the channel. Vehicles also had no revocation
information at all: they accepted messages from revoked senders until those senders' tickets
expired. This chapter describes the mechanisms added to address these problems. Each is an option
of the implementation and is measured on its own in the ablation study
(Section~\ref{sec:res-ablation}). Table~\ref{tab:profiles} groups them into the profiles compared in
Chapter~\ref{ch:results}.

\begin{table}[htbp]\centering\small
\caption{Mechanisms per profile ($\bullet$ = enabled)}\label{tab:profiles}
\resizebox{\textwidth}{!}{%
\begin{tabular}{@{}lcccc@{}}
\toprule
Mechanism & Phase-I & v2 & v3-fast & v3-private \\
\midrule
Ticket digests, compact wire format, handover hysteresis & & $\bullet$ & $\bullet$ & $\bullet$ \\
Fast quarantine over the RSU backbone & & $\bullet$ & $\bullet$ & $\bullet$ \\
TESLA with adaptive signing and key messages & & & $\bullet$ & $\bullet$ \\
ECQV implicit-certificate tickets & & & $\bullet$ & $\bullet$ \\
RSU clusters ($2\times2$) & & & $\bullet$ & $\bullet$ \\
Load-triggered batch verification & & & $\bullet$ & $\bullet$ \\
Dead-reckoning BSM triggering & & & $\bullet$ & $\bullet$ \\
Revocation filter in beacons & & & $\bullet$ & $\bullet$ \\
Trust-weighted reports & & & $\bullet$ & $\bullet$ \\
Ticket renewal; signal-strength check & & & $\bullet$ & $\bullet$ \\
BBS anonymous credentials & & & & $\bullet$ \\
Cooperative mix zones & & & & $\bullet$ \\
\bottomrule
\end{tabular}}
\end{table}

\section{Smaller Messages: Digests, Compact Encoding, Implicit Certificates}

A vehicle sends its full ticket only every tenth BSM, after hearing a new RSU, and at least every
2\,s; the other BSMs carry a 16-byte ticket digest that receivers resolve from a cache. The compact
wire format packs the kinematics into 16 bytes and uses 64-byte BIP-340 signatures. With ECQV
(Section~\ref{sec:tesla-ecqv}) the RSU issues an implicit certificate instead of signing the ticket, which saves the
explicit key and signature. The certificate request point is bound into the proof context, so nobody
can swap it on the air and receive a certificate for someone else's pseudonym.

\section{Cheaper Authentication: TESLA with Adaptive Signing}

With TESLA a BSM carries a 16-byte MAC under the current chain key instead of a signature. Keys are
disclosed two 25\,ms intervals later, in the vehicle's next BSM or in a small key message, so a
receiver can accept the BSM after about 50\,ms. A signed anchor BSM carries the chain commitment. To
avoid the TESLA delay when it brings nothing, the vehicle signs every BSM while its channel is quiet
and switches to TESLA when the channel busy ratio exceeds 0.30 (and back below 0.20, with 1\,s of
hysteresis). Urgent BSMs, after braking harder than 3\,m/s$^2$, are always signed so that neighbours
can act on them at once.

\section{Fewer Messages: Dead Reckoning}

Receivers predict a vehicle's position from its last BSM by dead reckoning. The vehicle sends a BSM
only when that prediction drifts by more than 1\,m, its speed changes by more than 0.5\,m/s or its
heading by more than $4^\circ$, and at least once a second. This follows the generation rules of the
European cooperative awareness service~\cite{etsi302637}. On straight roads at constant speed most
BSMs become unnecessary, while every manoeuvre is still reported immediately.

\section{Fewer Handovers: RSU Clusters and Ticket Renewal}

RSUs are grouped into clusters of $2\times2$, and a ticket is valid at every RSU of its cluster, so
moving between them needs no messages at all. Before a ticket expires, the vehicle renews it with one
signed request instead of a new proof.

\section{Load-Triggered Batch Verification}

Schnorr signatures can be verified in batches with one multi-scalar multiplication~\cite{bellare1998batch}.
The measured speed-up is \R{ct}{batch_speedup_8}$\times$ for 8 signatures and
\R{ct}{batch_speedup_32}$\times$ for 32. Waiting to fill a batch would add latency, so batching is
load-triggered: a signature that finds the processor idle is verified at once, and only signatures
that arrive while it is busy are verified together, up to 32, the moment it is free. If a batch fails,
bisection finds the invalid signatures.

\section{Revocation that Reaches Vehicles}

\textbf{Revocation filter.} Every fifth beacon (2\,Hz) carries a Bloom filter (7 hash functions, about 1\,\% false positives) of the pseudonyms
revoked or quarantined within the last ticket lifetime. Vehicles drop BSMs from pseudonyms in the
filter of the last beacon they verified. False positives are possible but rare, and they only drop
a message.

\textbf{Fast quarantine.} An RSU that reports a pseudonym also sends a signed QUARANTINE note over
the backbone to its neighbouring RSUs. Once notes from $K = 2$ distinct RSUs arrive, they block the
pseudonym locally until the ledger revocation lands, which takes milliseconds instead of a ledger
commit.

\section{Trust-Weighted Reports}\label{sec:trust}

A compromised RSU could report honest vehicles. With trust weighting, the TA investigates every
traced pseudonym. In the simulation the ground truth stands in for the investigation. Each RSU that
reported it gets a weight
\begin{equation}
  w = \min\!\left(1,\ \frac{2\,(\mathit{tp}+1)}{\mathit{tp}+\mathit{fp}+2}\right),
\end{equation}
where $\mathit{tp}$ and $\mathit{fp}$ count its confirmed and wrong reports. Revocation then needs a
report weight of at least $K$ instead of $K$ reporters. An honest RSU keeps weight 1; an RSU with more
wrong than right reports loses influence. Weights are stored on the ledger through
\texttt{SetReporterWeight}.

\section{Signal-Strength Check}

An RSU estimates the sender's distance from the received power using the scenario's path-loss model
with 4\,dB log-normal shadowing. It flags a BSM whose claimed position disagrees with that estimate
by more than four standard deviations and by at least 100\,m. A position-forging attacker cannot
move its own transmitter, so its claims stop matching its signal.

\section{The Privacy Profile: BBS Credentials and Mix Zones}

\textbf{BBS credentials.} The TA signs each vehicle's hidden secret, class and epoch once. A login
proves possession of that signature in zero knowledge and reveals the pseudonym
$\mathit{nym}_j = H(\mathit{epoch}, j)^{sk_v}$ for some $j < 32$. Logins are unlinkable even to
colluding RSUs, the TA no longer knows the pseudonyms in advance, and the identity escrow lists the
pseudonyms the credential can produce. The proof is about 70 times more expensive than a Schnorr proof
(Table~\ref{tab:compute}). Its nonce-independent part is therefore precomputed while the vehicle is
idle.

\textbf{Cooperative mix zones.} Changing a pseudonym alone does not stop a global eavesdropper, which
links the old and new pseudonyms by position and speed. Within 30\,m of a junction, a vehicle with at
least three neighbours triggers a coordinated change: every neighbour in the zone whose pseudonym is
old enough changes too, and each stays silent until it has left the zone and for a random
0.5--1.5\,s, at most 5\,s. The tracker attack of Section~\ref{sec:res-privacy} measures the effect.

\section{Emergency Vehicles}\label{sec:emergency}

Emergency vehicles need the opposite of anonymity on one point: their privileged status must be
verifiable by everyone, and must be impossible to fake. The extension builds this on the
authorization predicate (\ref{eq:access}).
\begin{enumerate}
  \item \textbf{Registration.} The TA registers an ambulance's pseudonyms with class
  \textit{emergency}. The class is in the Merkle leaf (or the BBS credential) and in the proof context,
  so a car cannot claim it.
  \item \textbf{Service.} Ambulances log in for the service \texttt{emergency-priority}, whose ledger
  policy admits only class \textit{emergency}. Their tickets therefore carry, signed by the RSU, the
  class and the service.
  \item \textbf{Siren alerts.} With the siren on, every BSM of an emergency vehicle is signed (never
  TESLA, so receivers need not wait), carries the full ticket and a flag \texttt{ev} inside the
  signed payload, and is never suppressed by dead reckoning. Vehicles process alerts from up to
  300\,m, beyond the 100\,m safety radius of normal BSMs.
  \item \textbf{Verification.} A receiver honours the flag only if the RSU-signed class of the
  sender's ticket is \textit{emergency}; otherwise it rejects the message as a fake. A signed fake is
  itself proof of the lie: no detector judgment is involved. The first RSU that receives one reports
  the pseudonym at once with the message hash as evidence, and quarantine and revocation follow.
  \item \textbf{Priority.} Emergency logins, renewals and handovers, and the RSUs' answers to
  them, are sent on the 802.11p access category AC\_VI instead of AC\_BE, so they do not wait
  behind the queue of ordinary logins on the radio. Safety messages keep the highest category,
  AC\_VO. The RSU also verifies them ahead of its processing queue: such a login waits only for the
  job in progress. Because the claim is checked only by the verification itself, an RSU grants
  priority to at most five claims per second, so false claims cannot crowd out everyone else.
\end{enumerate}
The fake-emergency attacker of the evaluation tries one emergency login, which the policy denies,
and then sets the flag on all of its BSMs.
